%% kapa: exam=1 round=28 year=2017 period=1 form=A subject=civil_law subject_ko=민법 seq=38 number=38
%% source: Q-Net 감정평가사 기출문제 (cst00309) artlSeq=5199603 file=extracted/1008/5199603/2017년 제28회 감정평가사 1차시험 문제지/1교시 A형.pdf
%% conversion: kordoc 4.24.51 -> tools/gen1.py (auto); review=proofread-sonnet; flags=none
\begin{question}{38}
근저당권에 관한 설명으로 옳은 것만을 모두 고른 것은? (다툼이 있으면 판례에 따름)
\begin{framedtext}
ㄱ. 피담보채무의 확정 전 채무자가 변경된 경우, 변경 후의 채무자에 대한 채권만이 당해 근저당권에 의하여 담보된다. \par
ㄴ. 근저당권의 존속기간이나 결산기의 정함이 없는 경우, 근저당권설정자는 근저당권자를 상대로 언제든지 해지의 의사표시를 함으로써 피담보채무를 확정시킬 수 있다. \par
ㄷ. 근저당권자가 피담보채무의 불이행을 이유로 경매신청을 한 경우, 경매신청시에 근저당권이 확정된다. \par
ㄹ. 선순위 근저당권의 확정된 피담보채권액이 채권최고액을 초과하는 경우, 후순위 근저당권자가 선순위 근저당권의 채권최고액을 변제하더라도 선순위 근저당권의 소멸을 청구할 수 없다.
\end{framedtext}
\begin{choices}
  \item ㄱ, ㄴ
  \item ㄴ, ㄷ
  \item ㄴ, ㄹ
  \item ㄱ, ㄷ, ㄹ
  \item ㄱ, ㄴ, ㄷ, ㄹ
\end{choices}
\end{question}
\answer{5}
